\scanpage{466}\mspage{378}
% archive sheet number: 465
\deleted{$M$ soit stable par $\varinjlim$ \uline{de type} $\Psi$, par $\varprojlim$ de type $\Phi$. Soit d'autre part $N$ une catégorie accessible, et considérons la \deleted{catégorie}}
\[
\underline{\mathrm{Hom}}_{\deleted{\unsure{\Phi}}\Psi}^{\deleted{\unsure{\Psi}}\Phi}(M^\circ,N)
\underset{\unsure{-|\to}}{\sim}
\left(\underline{\mathrm{Hom}}_{\Phi}^{\Psi}(M,N^\circ)\right)^\circ
\]
est accessible.

\uline{Dém} On procède comme pour 4.14.5, pour expliciter
\[
\underline{\mathrm{Hom}}_{\Phi}^{\Psi}(M,N^\circ)\xrightarrow{\approx}
\underline{\mathrm{Hom}}_{\Phi_\pi}^{\Psi}(C,N^\circ)
\]
d'où
\[
\underline{\mathrm{Hom}}_{\Psi}^{\Phi}(M^\circ,N)
\underset{\unsure{\leftrightarrow|\leftarrow|\to}}{\approx}
\underline{\mathrm{Hom}}_{\Psi}^{\Phi_\pi}(C^\circ,N),
\]
et on applique 4.14.2 \unsure{à} $C^\circ,\Psi,\Phi_\pi$.

\uline{Corollaire} 4.14.8. Soi\deleted{\unsure{t}}\unsure{ent} $M$ et $N$ \unsure{deux|des} catégories \deleted{accessibles, supposons \deleted{\unsure{par tt}} $M$ stable par petites $\varinjlim$. Alors la catégorie} \inserted{paratopos. Alors}
\[
M\boxtimes N\underset{\unsure{-|\to}}{\approx}\underline{\mathrm{Hom}}^!(M^\circ,N)
\qquad\text{(cf 4.2 et 4.3)}
\]
est un paratopos.

En effet, en vertu de 4.14.7 (\unsure{avec|cas} $\Phi=\Phi_{\unsure{M|1|0}}$, $\Psi=\varnothing$) c'est une catégorie accessible.
